\documentclass[12pt,a4paper,bibliography=totocnumbered,listof=totocnumbered]{scrartcl}
\usepackage[ngerman]{babel}
\usepackage[utf8]{inputenc}
\usepackage{amsmath}
\usepackage{amsfonts}
\usepackage{amssymb}
\usepackage{graphicx}
\usepackage{fancyhdr}
\usepackage{tabularx}
\usepackage{geometry}
\usepackage{setspace}
\usepackage[right]{eurosym}
\usepackage[printonlyused]{acronym}
\usepackage{subfig}
\usepackage{floatflt}
\usepackage[usenames,dvipsnames]{color}
\usepackage{colortbl}
\usepackage{paralist}
\usepackage{array}
\usepackage{titlesec}
\usepackage{parskip}
\usepackage[right]{eurosym}
\usepackage{picins}
\usepackage[subfigure,titles]{tocloft}
\usepackage[pdfpagelabels=true]{hyperref}

\usepackage{listings}
\lstset{basicstyle=\footnotesize, captionpos=b, breaklines=true, showstringspaces=false, tabsize=2, frame=lines, numbers=left, numberstyle=\tiny, xleftmargin=2em, framexleftmargin=2em}
\makeatletter
\def\l@lstlisting#1#2{\@dottedtocline{1}{0em}{1em}{\hspace{1,5em} Lst. #1}{#2}}
\makeatother

\geometry{a4paper, top=27mm, left=30mm, right=20mm, bottom=35mm, headsep=10mm, footskip=12mm}

\hypersetup{unicode=false, pdftoolbar=true, pdfmenubar=true, pdffitwindow=false, pdfstartview={FitH},
	pdftitle={Bachelorarbeit},
	pdfauthor={Mathias Möller},
	pdfsubject={Bachelorarbeit},
	pdfcreator={\LaTeX\ with package \flqq hyperref\frqq},
	pdfproducer={pdfTeX \the\pdftexversion.\pdftexrevision},
	pdfkeywords={Bachlorarbeit},
	pdfnewwindow=true,
	colorlinks=true,linkcolor=black,citecolor=black,filecolor=magenta,urlcolor=black}
\pdfinfo{/CreationDate (D:20110620133321)}

\begin{document}

\titlespacing{\section}{0pt}{12pt plus 4pt minus 2pt}{-6pt plus 2pt minus 2pt}

% Kopf- und Fusszeile
\renewcommand{\sectionmark}[1]{\markright{#1}}
\renewcommand{\leftmark}{\rightmark}
\pagestyle{fancy}
\lhead{}
\chead{}
\rhead{\thesection\space\contentsname}
\lfoot{Entwurf und Implementierung eines Chrome Plugins zur automatischen\newline Anreicherung von Webseiten mit kulturellen Inhalten}
\cfoot{}
\rfoot{\ \linebreak Seite \thepage}
\renewcommand{\headrulewidth}{0.4pt}
\renewcommand{\footrulewidth}{0.4pt}

% Vorspann
\renewcommand{\thesection}{\Roman{section}}
\renewcommand{\theHsection}{\Roman{section}}
\pagenumbering{Roman}

% ----------------------------------------------------------------------------------------------------------
% Titelseite
% ----------------------------------------------------------------------------------------------------------
\thispagestyle{empty}
\begin{center}
	\includegraphics[scale=0.3]{Bilder/uni_passau.png}\\
	\vspace*{2cm}
	\Large
	\textbf{Fakultät für}\\
	\textbf{Informatik und Mathematik}\\
	\vspace*{2cm}
	\Huge
	\textbf{Bachelorarbeit}\\
	\vspace*{0.5cm}
	\large
	über das Thema\\
	\vspace*{1cm}
	\textbf{Entwurf und Implementierung eines Chrome Plugins zur automatischen Anreicherung von Webseiten mit kulturellen Inhalten}\\
	\vspace*{2cm}
	
	\vfill
	\normalsize
	\newcolumntype{x}[1]{>{\raggedleft\arraybackslash\hspace{0pt}}p{#1}}
	\begin{tabular}{x{6cm}p{7.5cm}}
		\rule{0mm}{5ex}\textbf{Autor:} & Mathias Möller\newline moellerm@fim.uni-passau.de \\ 
		\rule{0mm}{5ex}\textbf{Prüfer:} & Prof. Granitzer \\ 
		\rule{0mm}{5ex}\textbf{Abgabedatum:} & 24.08.2015 \\ 
	\end{tabular} 
\end{center}
\pagebreak

% ----------------------------------------------------------------------------------------------------------
% Abstract
% ----------------------------------------------------------------------------------------------------------
\setcounter{page}{1}
\onehalfspacing
\titlespacing{\section}{0pt}{12pt plus 4pt minus 2pt}{2pt plus 2pt minus 2pt}
\rhead{KURZFASSUNG}
\section{Kurzfassung}
Lorem ipsum dolor sit amet, consetetur sadipscing elitr, sed diam nonumy eirmod tempor invidunt ut labore et dolore magna aliquyam erat, sed diam voluptua. At vero eos et accusam et justo duo dolores et ea rebum. Stet clita kasd gubergren, no sea takimata sanctus est Lorem ipsum dolor sit amet. Lorem ipsum dolor sit amet, consetetur sadipscing elitr, sed diam nonumy eirmod tempor invidunt ut labore et dolore magna aliquyam erat, sed diam voluptua. At vero eos et accusam et justo duo dolores et ea rebum. Stet clita kasd gubergren, no sea takimata sanctus est Lorem ipsum dolor sit amet.

\vspace{-1,2em}
\titlespacing{\section}{0pt}{12pt plus 4pt minus 2pt}{-6pt plus 2pt minus 2pt}
\section*{Abstract}
Das ganze auf Englisch.

\pagebreak

% ----------------------------------------------------------------------------------------------------------
% Verzeichnisse
% ----------------------------------------------------------------------------------------------------------
% TODO Typ vor Nummer
\renewcommand{\cfttabpresnum}{Tab. }
\renewcommand{\cftfigpresnum}{Abb. }
\settowidth{\cfttabnumwidth}{Abb. 10\quad}
\settowidth{\cftfignumwidth}{Abb. 10\quad}

\titlespacing{\section}{0pt}{12pt plus 4pt minus 2pt}{2pt plus 2pt minus 2pt}
\singlespacing
\rhead{INHALTSVERZEICHNIS}
\renewcommand{\contentsname}{II Inhaltsverzeichnis}
\phantomsection
\addcontentsline{toc}{section}{\texorpdfstring{II \hspace{0.35em}Inhaltsverzeichnis}{Inhaltsverzeichnis}}
\addtocounter{section}{1}
\tableofcontents
\pagebreak
\rhead{VERZEICHNISSE}
\listoffigures
\pagebreak
\listoftables
%\pagebreak
\renewcommand{\lstlistlistingname}{Listing-Verzeichnis}
{\labelsep2cm\lstlistoflistings}
\pagebreak

% ----------------------------------------------------------------------------------------------------------
% Abkürzungen
% ----------------------------------------------------------------------------------------------------------
\section{Abk\"urzungsverzeichnis}
\begin{acronym}[JITIR] % längste Abkürzung steht in eckigen Klammern
	\setlength{\itemsep}{-\parsep} % geringerer Zeilenabstand
	\acro{JITIR}{Just-in-time information retrieval}
\end{acronym}
\newpage

% ----------------------------------------------------------------------------------------------------------
% Inhalt
% ----------------------------------------------------------------------------------------------------------
% Abstände Überschrift
\titlespacing{\section}{0pt}{12pt plus 4pt minus 2pt}{-6pt plus 2pt minus 2pt}
\titlespacing{\subsection}{0pt}{12pt plus 4pt minus 2pt}{-6pt plus 2pt minus 2pt}
\titlespacing{\subsubsection}{0pt}{12pt plus 4pt minus 2pt}{-6pt plus 2pt minus 2pt}

% Kopfzeile
\renewcommand{\sectionmark}[1]{\markright{#1}}
\renewcommand{\subsectionmark}[1]{}
\renewcommand{\subsubsectionmark}[1]{}
\lhead{Kapitel \thesection}
\rhead{\rightmark}

\onehalfspacing
\renewcommand{\thesection}{\arabic{section}}
\renewcommand{\theHsection}{\arabic{section}}
\setcounter{section}{0}
\pagenumbering{arabic}
\setcounter{page}{1}

\section{Einleitung und Motivation}
Suchmaschinen gehören heutzutage zu den meistbesuchten Seiten im Internet. Google\footnote{https://www.google.com/} ist mit 2.000.000.000.000 Suchanfragen pro Jahr (TODO: Quelle) nicht nur die beliebteste Suchmaschine, sondern auch die meist besuchte Seite im Internet. Auch die chinesische Suchmaschine Baidu\footnote{http://www.baidu.com/} ist in den oberen fünf Plätzen vertreten \cite{alexa}. 
\begin{itemize}
	\item Auf XX\% der Webseiten wird über eine Suchmaschine zugegriffen
	\item Menschen holen Informationen über Suchmaschinen, nicht mehr aus Bibliotheken oder Lexika 
	\item Woher Quellen?
\end{itemize}
Um so mehr sie an Bedeutung gewinnen, umso größer ist der Bedarf an bessern Suchmaschinen \cite{lawrence2000context}.

Doch der Funktionsumfang klassischer Suchmaschinen ist begrenzt. Sie können den Kontext des Benutzers in ihren Queries? nicht miteinbeziehen. Weiterhin verlangen sie vom Nutzer, seine gegenwärtige Arbeit einzustellen, die Webseite der Suchmaschine aufzurufen und eine Suchanfrage zu formulieren. So wird die Konzentration auf den eigentlichen Arbeitsschritt gestört. Im Schnitt sind nur 67\% der Suchanfragen an Google erfolgreich (TODO: Quelle). Das heißt in circa einem Drittel der Fälle bleibt der Aufruf der Suchmaschine ohne die erwarteten Erfolge. Zu der verlorenen Zeit addiert sich auch noch die Zeit, die der Nutzer braucht um sich seine Arbeit wieder ins Kurzzeitgedächtnis zu holen.

Um diesen Problemen entgegen zu wirken wurden neue Wege zur Informationsgewinnung entwickelt. Eine Möglichkeit sind so genannte Just-in-time Information Retrieval Agents (JITIR-Agents) \cite{rhodes2000just}. Sie beobachten im Hintergrund den Kontext des Benutzers und versuchen aus den so erhaltenen Informationen eine Suchanfrage an eine Datenbank oder ein Recommender-System zu schicken. Die gewonnenen Informationen werden dann möglichst unaufdringlich dem Benutzer angezeigt. Er kann sich nun entscheiden diese Informationen genauer zu betrachten oder mit seiner Arbeit fortzufahren. Die kognitive Belastung bleibt hierbei sehr gering. JITIR-Agents reduzieren auf diese Weise enorm den Aufwand Informationen zu finden \cite{rhodes2000just}.

Durch ihre Funktionsweise sind sie jedoch nicht so exakt wie klassische Suchmaschinen, da sie nur ``erraten'' können, was den Benutzer gerade interessiert. Wenn ein Benutzer einer genaue Vorstellung von der Suchanfrage oder den Ergebnissen hat, haben klassische Suchmaschinen Vorteile gegenüber den JITIR-Agents \cite{rhodes2000just}.

Ein weiteres Problem des Internets ist es, dass die Standards, die die simple und leicht skalierbare Architektur des Netzwerks ermöglichen, den Einsatz von fortgeschritteneren Hypermedia-Technologien verhindern \cite{bouvin1999unifying}. Eine Möglichkeit, das Internet um Funktionen zu erweitern ist Web Augmentation (WA). Diaz \cite{diaz2012understanding} beschreibt WA als den Versuch, statt eine neue Technologie zu entwickeln, neue Funktionalität auf eine gerenderte Webseite zu setzen:
\begin{quote}
In some sense, WA is to the Web what Augmented Reality is to the physical world: layering relevant content/layout/navigation over the existing Web to customize the user experience. \cite{diaz2012understanding}
\end{quote}
Solche WA-Tools interagieren mit einem Web Server, Http-Proxy oder dem Browser der Nutzers, um so Inhalte oder Navigationselemente direkt in die angezeigte Webseite einzufügen. Auf diese Weise erlauben sie es, die limitierten Möglichkeiten des World Wide Webs zu Gunsten des Nutzers anzureichern (QUelle: Kennth anderson oder Bouvin).

Ein Beispiel für ein solches WA-Tool sowie für ein JITIR System ist das EEXCESS Chrome Plugin\footnote{http://eexcess.eu/results/chrome-extension/}. Es analysiert die aktuell besuchten Seiten des Nutzers und schlägt ihm auf Grund der erlangten Daten am Rand des Browser-Fensters weiterführende Quellen aus der Europeana-Datenbank\footnote{http://www.europeana.eu/portal/} vor.

Ziel dieser Bachelorarbeit ist es, auf Basis des EEXCESS Plugins ein Chrome Plugin zu entwerfen und zu implementieren, welche eine Webseite in einzelne Paragraphen aufteilt und zu jedem dieser Paragraphen kulturelle Inhalte der Europeana-Datenbank vorschlägt. Dabei soll das Design des Plugins helfen, die bei der in Kapitel 5 erläuterten Evaluierung des EEXCESS Plugins aufgetauchten Probleme zu beheben und Schwachstellen zu verbessern.

Im zweiten Teil dieser Arbeit wird auf vergleichbare Technologien eingegangen und ihre Gemeinsamkeiten und Unterschiede zu diesem Projekt. Anschließend wird das Konzept und die Implementierung des Projekts beschrieben und das Ergebnis evaluiert. Der sechste Teil der Ausarbeitung widmet sich Möglichkeiten für zukünftige Projekte und Verbesserungen. Zuletzt wird ein Fazit gezogen und die Ergebnisse der Arbeit analysiert.
\section{Konzept}
 \subsection{Warum kein Proactiver JITIR-Agent?}
 	- > API Limitierung und decrease cognitive load 
	- > Benutzer entscheidet ob er weitere Informationen erhalten möchte
	- > Benutzer kann Suchanfrage erst anpassen (Nachteil von Margin Notes (Paper 4)
 \subsection{Anzeige der Ergebnisse}
 \subsection{Erklärung der Such-Anfragen Generierung}
 \subsection{Anpassen der Suchanfrage durch den Nutzer}
 \subsection{Verbesserung der Suchanfrage z.B. durch maschinelles Lernen}

\section{Implementierung}
 \subsection{Verwendung von AngularJS für alle Komponenten des Plugins}
 \subsection{Bau der GUIs}
		- > Darf den Benutzer nicht zu sehr ablenken
		- > Ergebnisse müssen in der Nähe ihrer „Quelle“ angezeigt werden (proximity compatibility pricinple)
		- > Benutzer muss klar zwischen Webseite und Augmentation unterscheiden können
		- > buntes, auffälliges Design
		- > Ramping interface: Mehr Benutzerinteraktionen führen zu mehr angezeigten Informationen (Erklärung der Stages) 
 \subsection{Einbindung der REST-Services}

\section{Evaluierung?}
Aufgaben die Benutzer mit EEXCESS Lösen mussten müssen sie jetzt mit Redesign lösen. Vergleich der Ergebnisse?

\section{Future Work}
 \subsection{Alternative Algorithmen zur Textanalyse}
 \subsection{Implementierung einer automatischen Suchanfragen-Verbesserung durch maschinelles Lernen}
 	\begin{itemize}
 		\item Mehr kontextuelle Informationen miteinbeziehen
 		\item Such-Profil des Nutzers erstellen

 	\end{itemize}
 \subsection{Verbesserung der Ergebnis-Güte}
	\begin{itemize}
		\item durch Query Expansion
 		\item durch Filtern der Ergebnisse (mehr Präzision da Ausbeute bei JITIR nicht so relevant) Clustering
 		\item relevance feedback
	\end{itemize}
\subsection{Anbindung weiterer Quellen}
	Automatische Auswahl der richtigen QUelle laut selecting task relevant sources
 \subsection{Anpassung der Anwendung auf mobile Nutzung}

 \section{Conclusion}
 \begin{itemize}
 		\item Steigerung der Effektivität und Produktivität von wissenschaftlichem Arbeiten
 		\item Starke Effektivitätssteigerung durch Punkte aus Future Work möglich?
 		\item information management assistants embody a vision of a future in which users hardly ever form a query to request information. when an information need arises, a system like watson has already anticipated it and provided relevant information to the user before she is even able to ask for it (budzik, watson)
 		\item survey suggested that users are relatively dissatisfied with the results of their searching experience. THis makes concrete the claim that systems designed to help useres in their information seeking tasks are needed in the world (budzik, watson)

 		JITIR's can be thought of as automatic ``query free'' search engines (rhodes, using physical context)
 	\end{itemize}

% ----------------------------------------------------------------------------------------------------------
% Literatur
% ----------------------------------------------------------------------------------------------------------
\renewcommand\refname{Quellenverzeichnis}
\bibliographystyle{myalpha}
\bibliography{bibo}
\pagebreak

% ----------------------------------------------------------------------------------------------------------
% Anhang
% ----------------------------------------------------------------------------------------------------------
\pagenumbering{Roman}
\setcounter{page}{1}
\lhead{Anhang \thesection}

\begin{appendix}
\section*{Anhang}
\phantomsection
\addcontentsline{toc}{section}{Anhang}
\addtocontents{toc}{\vspace{-0.5em}}


\subsection*{GUI Screenshots}
\label{app:screenshot}
Unterkategorie, die nicht im Inhaltsverzeichnis auftaucht.

\end{appendix}


\newpage
\thispagestyle{empty}
\begin{center}
	\vspace*{5em}
	\huge\textbf{Erklärung}\\
\end{center}
\vspace{2em}
Hiermit versichere ich, dass ich meine Abschlussarbeit selbständig verfasst und keine anderen als die angegebenen Quellen und Hilfsmittel benutzt habe.

\vspace{4em}
\begin{minipage}{\linewidth}
	\begin{tabular}{p{15em}p{15em}}
		Datum: &  .......................................................\\
		& \centering (Unterschrift)\\
	\end{tabular}
\end{minipage}

\end{document}
